% Generated draft: requires mathematical and source review.
\documentclass[11pt]{article}
\usepackage[margin=25mm]{geometry}
\usepackage{fontspec}
\setmainfont{Latin Modern Roman}
\usepackage{amsmath,amssymb,mathtools,unicode-math}
\setmathfont{Latin Modern Math}
\usepackage{enumitem,needspace,xurl,hyperref}
\setlength{\parindent}{0pt}
\setlength{\parskip}{5pt}
\setlength{\emergencystretch}{3em}
\newcommand{\N}{\mathbb{N}}
\newcommand{\Z}{\mathbb{Z}}
\newcommand{\Q}{\mathbb{Q}}
\newcommand{\R}{\mathbb{R}}
\newcommand{\C}{\mathbb{C}}
\newcommand{\F}{\mathbb{F}}
\newcommand{\A}{\mathbb{A}}
\newcommand{\PP}{\mathbb{P}}
\newcommand{\E}{\mathbb{E}}
\newcommand{\eps}{\varepsilon}
\newcommand{\dd}{\,\mathrm{d}}
\begin{document}
\section*{UnsolvedMath ID: 30001987}
\textbf{Source problem code:} catalogue entry
\subsection{Definitions and mathematical statement}
A differential field with one derivation is a field \(K\) equipped with a map \(\delta_K:K\to K\) satisfying
\[
\delta_K(a+b)=\delta_K(a)+\delta_K(b),\qquad
\delta_K(ab)=a\delta_K(b)+b\delta_K(a)
\quad(a,b\in K).
\]
For differential fields \((K,\delta_K)\) and \((L,\delta_L)\), a differential-field embedding is an injective field homomorphism \(g:K\to L\) such that
\[
g(\delta_K(a))=\delta_L(g(a))\qquad(a\in K).
\]
These definitions describe the single-derivation setting; the supplied fragment does not specify the number of derivations, the characteristic, or the hypotheses of Harrington's construction.

Given presentations in which elements have natural-number codes, an embedding is computable when its action on those codes is computable. Its image is
\[
g(K)=\{g(a):a\in K\}\subseteq L.
\]
For a specified coding \(\iota:L\to\N\), the degree of this image means
\[
\deg_T\bigl(\{\iota(g(a)):a\in K\}\bigr),
\]
where \(A\leq_T B\) means that membership in \(A\subseteq\N\) is decidable by an algorithm with oracle \(B\), and \(\deg_T(A)\) is the equivalence class of \(A\) under mutual Turing reducibility. The relevant presentations and coding must be specified to interpret this expression.

For comparison, a usual splitting-set convention for a computably presented field \(F\), with an effective coding of \(F[X]\), is
\[
\operatorname{Split}(F)=
\{\ulcorner f\urcorner: f\in F[X],\ \deg f\geq2,
\ \exists u,v\in F[X]\ (\deg u,\deg v\geq1\ \land\ f=uv)\}.
\]
This convention is explanatory; the supplied record does not define its intended splitting set or identify the differential property that should replace polynomial reducibility.

\textbf{Source question, with unresolved parameters.} For the computable differential-field embedding \(g\) attributed in the record to Harrington, determine the Turing degree of \(g(K)\), and identify an appropriate differential analogue of the splitting set. The record does not specify \(K\), the target field, their presentations, the construction or choices defining \(g\), or the criterion that makes a differential analogue appropriate. It also does not state whether a single degree or a classification over a family of embeddings is sought. Consequently, this is a faithful statement of the extracted question with explicit gaps, not a fully specified mathematical problem.
\subsection{Short English statement}
The question asks how computationally complicated the image of Harrington's computable differential-field embedding is, and what should play the role of a splitting set for differential fields. The supplied fragment omits the embedding's construction, field hypotheses, presentations, and the intended differential analogue, so these parts require source clarification before a precise target can be fixed.
\subsection{Sources}
\begin{itemize}
\item Local UnsolvedMath registry, permanent ID \(30001987\), title \emph{Turing Degrees of Differential-Field Embedding Images}, source problem code \texttt{catalogue entry}, in the Oberwolfach Reports Open Problems collection. The supplied record gives proposed year 2012 and no named proposer; it attributes the embedding to Harrington.
\item Supplied report citation: \emph{Computability Theory} (2012), Oberwolfach Reports, pp.\ 397--455, DOI \texttt{10.4171/owr/2012/08}. Source URL and supplied literature-source link labelled ``original Oberwolfach report'': \url{https://doi.org/10.4171/owr/2012/08}.
\item Literature assessment dated exactly \texttt{2026-08-21}: supplied status \texttt{open}. The supplied assessment describes a historical fragment about reducibility of polynomial representations, says that it does not state a complete decidable theorem, and describes the underlying question as unresolved as extracted. This is supplied dated curation, not an independently verified literature assessment; the record requests independent verification of time-sensitive status and cited claims.
\item The registry supplies statement status \texttt{corrected\_verified} and reports that the mathematical content was checked against the official report, with notation, context, and question form normalized without changing the source-level claim. This draft preserves that verification claim as source metadata and does not independently confirm it.
\end{itemize}
\end{document}
